\section{Thiết kế Domain}

\subsection{Mô hình phân cấp tài nguyên}

Mô hình dữ liệu nghiệp vụ phản ánh phân cấp vật lý của hệ thống như
Hình~\ref{fig:resource-hierarchy}. Quan hệ Cluster--Node và Node--Runtime là
một--nhiều; Node--Agent là một--một nên Agent không cần entity riêng mà được
biểu diễn qua các thuộc tính của Node (\texttt{agent\_version},
\texttt{last\_heartbeat\_at}).

\begin{figure}[H]
\centering
\begin{tikzpicture}[node distance=0.45cm and 0.5cm]
  \node[groupbox] (cloud) {Cloud};
  \node[groupbox, below=of cloud] (cluster) {Edge Cluster};
  \node[groupbox, below=of cluster] (node) {Edge Node};
  \node[groupbox, below=of node] (agent) {Edge Agent};
  \node[groupbox, below=of agent] (runtime) {AI Runtime};
  \draw[flowarrow] (cloud) -- node[right, font=\scriptsize]{1 : N} (cluster);
  \draw[flowarrow] (cluster) -- node[right, font=\scriptsize]{1 : N} (node);
  \draw[flowarrow] (node) -- node[right, font=\scriptsize]{1 : 1} (agent);
  \draw[flowarrow] (agent) -- node[right, font=\scriptsize]{1 : N} (runtime);
\end{tikzpicture}
\caption{Phân cấp tài nguyên của nền tảng}
\label{fig:resource-hierarchy}
\end{figure}

\subsection{Cloud domain}

Bảng~\ref{tab:cloud-entities} tóm tắt các entity của Cloud. Entity được hiện thực
bằng \texttt{dataclass}, không phụ thuộc ORM; repository chuyển đổi giữa entity và
bản ghi database.

\begin{table}[H]
\centering
\caption{Các entity chính của Cloud domain}
\label{tab:cloud-entities}
\begin{tabularx}{\textwidth}{L{2.4cm}L{5.2cm}Y}
\toprule
\textbf{Entity} & \textbf{Thuộc tính chính} & \textbf{Quy tắc nghiệp vụ}\\
\midrule
Cluster & id, name, status, location, metadata & Tên duy nhất trên toàn hệ thống.\\
Node & id, cluster\_id, name, hostname, ip, status, agent\_version, cpu\_cores, memory\_total\_mb, gpu\_name, last\_heartbeat\_at & Tên duy nhất trong cluster; thông tin phần cứng do Agent tự báo cáo; OFFLINE khi quá hạn heartbeat.\\
Runtime & id, node\_id, name, type, status, image, model\_id, config, environment, resource\_limits, container\_id & Trạng thái do Edge báo cáo; bị xóa khi Edge xác nhận gỡ bỏ.\\
Deployment & id, cluster\_id, node\_id, runtime\_id, spec, status, last\_command\_id, message & Desired state của một runtime trên một node; chỉ chấp nhận kết quả của lệnh mới nhất.\\
\bottomrule
\end{tabularx}
\end{table}

\textbf{Deployment là aggregate trung tâm.} Mỗi lần phát lệnh, phương thức
\texttt{dispatched(command\_id)} ghi lại \texttt{last\_command\_id} và chuyển trạng
thái. Khi Edge báo kết quả, \texttt{apply\_result(command\_id, status)} chỉ áp dụng
nếu \texttt{command\_id} trùng lệnh mới nhất; kết quả của lệnh cũ đến muộn bị bỏ
qua. Đây là cơ chế chống ghi đè ngược trạng thái khi mạng giao lại thông điệp
không theo thứ tự. Deployment ở trạng thái REMOVED là trạng thái kết thúc và
không nhận lệnh mới.

\begin{figure}[H]
\centering
\resizebox{\textwidth}{!}{%
\begin{tikzpicture}[font=\small,
  st/.style={draw, rounded corners, align=center, minimum height=0.8cm, minimum width=2.9cm, fill=blue!5},
  lbl/.style={font=\scriptsize, fill=white, inner sep=1.5pt}]
  \node[st] (pending) at (0,0) {PENDING};
  \node[st] (prog) at (4.2,0) {IN\_PROGRESS};
  \node[st] (ok) at (8.6,1.4) {SUCCEEDED};
  \node[st] (fail) at (8.6,-1.4) {FAILED};
  \node[st] (removing) at (13.0,0) {REMOVING};
  \node[st, fill=gray!15] (removed) at (17.2,0) {REMOVED};
  \draw[flowarrow] (pending) -- node[lbl, above=2pt]{dispatch} (prog);
  \draw[flowarrow] (prog.north east) -- node[lbl, above left]{kết quả OK} (ok.west);
  \draw[flowarrow] (prog.south east) -- node[lbl, below left]{kết quả lỗi} (fail.west);
  \draw[flowarrow] (ok.east) -- node[lbl, above right]{remove} (removing.north west);
  \draw[flowarrow] (fail.east) -- node[lbl, below right]{remove} (removing.south west);
  \draw[flowarrow] (removing) -- node[lbl, above=2pt]{Edge xác nhận} (removed);
  \draw[flowarrow, dashed] (ok.north) -- ++(0,0.7) -| node[lbl, pos=0.25, above]{update / start / stop / restart} (prog.north);
\end{tikzpicture}}
\caption{Máy trạng thái của Deployment}
\label{fig:deployment-state-machine}
\end{figure}

Value object gồm \texttt{ModelRef} (repo\_id, filename, revision),
\texttt{ModelLocation} (URL tải đã ghim theo commit SHA) và
\texttt{NodeHeartbeat} (các thông tin Node trích từ bản tin NodeState).

\subsection{Edge domain}

Phía Edge có entity \texttt{RuntimeInstance} gồm \texttt{RuntimeSpec} mong muốn,
\texttt{desired\_running}, trạng thái hiện tại, \texttt{container\_id},
\texttt{restart\_count}, số lần khởi chạy thất bại liên tiếp và thông điệp lỗi gần
nhất. \texttt{DesiredStateRepository} lưu danh sách RuntimeInstance xuống tệp JSON
trong \texttt{state\_dir} để Agent khôi phục sau khi khởi động lại.

Vòng đời runtime dùng chung một enum giữa hai phía (Bảng~\ref{tab:runtime-states}).

\begin{table}[H]
\centering
\caption{Trạng thái của AI Runtime}
\label{tab:runtime-states}
\begin{tabularx}{\textwidth}{L{3cm}Y}
\toprule
\textbf{Trạng thái} & \textbf{Ý nghĩa}\\
\midrule
CREATING & Agent đang tải model và pull image.\\
STARTING & Container đang được khởi chạy.\\
RUNNING & Container chạy bình thường.\\
DEGRADED & Container chạy nhưng hiệu năng suy giảm.\\
STOPPING / STOPPED & Đang dừng / đã dừng theo yêu cầu.\\
RESTARTING & Đang khởi động lại.\\
FAILED & Khởi chạy thất bại hoặc vượt số lần thử tối đa.\\
UNKNOWN & Chưa xác định (ví dụ ngay sau khi Agent khởi động lại).\\
\bottomrule
\end{tabularx}
\end{table}

\subsection{Hợp đồng dùng chung Cloud--Edge}

Gói \texttt{src/shared/contracts} định nghĩa bằng Pydantic toàn bộ thông điệp
trao đổi: \texttt{Command} và \texttt{RuntimeSpec} (Cloud$\rightarrow$Edge);
\texttt{NodeState}, \texttt{RuntimeState}, \texttt{SystemTelemetry},
\texttt{GpuTelemetry}, \texttt{RuntimeTelemetry}, \texttt{EdgeEvent},
\texttt{DeploymentEvent} (Edge$\rightarrow$Cloud) cùng các enum trạng thái. Giá
trị enum trùng khớp với kiểu enum trong PostgreSQL. Vì hai phía import cùng một
mã nguồn, thông điệp sai schema bị từ chối ngay tại nơi nhận thay vì làm hỏng dữ
liệu.

\texttt{RuntimeSpec} là desired state của một runtime:
\texttt{runtime\_id}, \texttt{name}, \texttt{type} (VISION, SAFETY, ANOMALY, OCR,
CUSTOM), \texttt{image\_uri}, \texttt{image\_tag}, \texttt{model\_id},
\texttt{model\_uri}, \texttt{config}, \texttt{environment} và
\texttt{resource\_limits}. Deploy và update là idempotent theo \texttt{runtime\_id}.
